\section{Models}
\label{section:models}

Our core model is \textbf{masked PPO}, implemented with \texttt{MaskablePPO} from
\texttt{sb3-contrib} (Stable-Baselines3 ecosystem), with \texttt{RecurrentPPO} available as an optional variant.
Using an established SB3 implementation was an intentional choice: it gives us a reliable optimization pipeline,
well-tested training code, and straightforward integration with our Gymnasium-compatible environment.

This choice is also motivated by the structure of Pokemon battles: the state is high-dimensional,
transitions are stochastic, and the action space is discrete with many context-dependent invalid actions.
In this regime, an on-policy actor-critic approach with conservative policy updates is generally more stable
than value-based baselines, and action masking improves sample efficiency by restricting exploration to legal
decisions only.

Operationally, at each turn the policy receives the encoded state vector \(\mathbf{s}_t\) from
Section~\ref{section:environment}, then outputs a categorical distribution over actions.
The environment mask is applied so only legal actions remain selectable; the policy samples (or chooses)
an action from that masked distribution. In parallel, the critic estimates \(V_\phi(\mathbf{s}_t)\), used
to compute an advantage estimate \(\hat{A}_t\) for policy optimization.
This actor-critic decomposition is useful here because it separates action selection from state evaluation:
the actor improves tactical choice quality, while the critic provides a lower-variance learning signal.

PPO updates the policy using the probability ratio
\(r_t(\theta)=\pi_\theta(a_t\mid s_t)/\pi_{\theta_{\text{old}}}(a_t\mid s_t)\), with clipped objective:
\[
L^{\text{clip}}(\theta)=\mathbb{E}_t\left[\min\left(r_t(\theta)\hat{A}_t,\;\mathrm{clip}(r_t(\theta),1-\epsilon,1+\epsilon)\hat{A}_t\right)\right].
\]
The clipping term (cut-off) prevents overly large policy updates in one step.
The final training loss also includes a value-function term and an entropy bonus to keep exploration.

In practice, the clipped ratio term is critical: without it, policy updates can become too aggressive,
especially in noisy environments where short-term returns fluctuate strongly from one rollout to another.
The clipped objective keeps updates within a trust region around the previous policy, which tends to produce
more stable training dynamics.

For completeness, we also implemented additional approaches in the same training framework
(notably \texttt{DQN}, and the recurrent PPO variant). However, these alternatives were not evaluated as
extensively as the default masked PPO setup, so our main analysis and conclusions focus on the SB3 masked PPO model.

In short, the policy learns to assign probability mass to strong legal actions while remaining conservative
between updates, which is particularly important in noisy, partially observable battles.
Implementation details (network architecture and hyperparameters) are reported in Appendix~\ref{section:appendix}.